\section{Finite normal pairs from the common construction}
\label{sec:pairs}
The complete ordered-pair encoder and its recognition, rigidity, extension,
color-kernel and chamber-completion proofs now belong to the common
paper~\cite{Recognition}. They are the constructions of D\'iaz Ramos,
Molinier and Viruel~\cite[Theorem 3.1 and Theorem A]{DMV}, with all
hypotheses and exceptional parameters retained. The present application
records both full source statements before treating the additional
recursive family and completed hexagon.

\begin{theorem}\label{R.FinitePairs}
Every finite normal pair $H\mathrel{\trianglelefteq}G$ is isomorphic to
$(C(\Gamma),K(\Gamma))$ for a finite connected incidence geometry $\Gamma$.
For $G\ne1$ the construction below has rank $[G:H]+3$; for $G=1$ it has rank
one. No fixed rank-two assertion is made.
\end{theorem}
\begin{proof}
Apply common Theorem UR.FinitePairs to the finite pair $(G,H)$.
It includes the explicit one-object geometry for $G=1$; otherwise its
ordered-pair graph is connected, triangle-free and has minimum degree two,
so the common chamber theorem gives the stated connected geometry and rank.
\end{proof}

\begin{theorem}[Ordered-pair encoding]\label{R.PairEncoding}
Let $G\ne1$ be finite and $H\mathrel{\trianglelefteq}G$. There is a connected
properly colored triangle-free graph $X(G,H)$ of minimum degree two with
\[
 (C(X(G,H)),K(X(G,H)))\cong(G,H).
\]
For $q=|G|$ it has $(q^3+9q^2-8q)/2$ vertices and $[G:H]+3$ colors.
\end{theorem}
\begin{proof}
Apply common Theorem UR.PairEncoding with $q=|G|\ge2$.
Its full ordered edge list, intrinsic degree recognition and unique
extension of right translations prove the exact pair, graph hypotheses,
vertex count and color count, including $q=2$.
\end{proof}

Normality is necessary because the type-preserving group is the kernel
of the palette action. In particular $G/H$ must embed in the symmetric
group on that palette; rank two cannot realize arbitrary indices.
The complete arbitrary-input cleaning/conversion theorem is also in the
common paper. It preserves components and does not silently make a
disconnected input connected.
